% One-time imported draft. Editable native source; no Markdown regeneration.
\section{The complete neighbor-count
rule}\label{008-life-generations-x-rule-the-complete-neighbor-count-rule}

{\def\LTcaptype{table} % do not increment counter
\begin{longtable}[]{@{}
  >{\RaggedRight\arraybackslash}p{(\linewidth - 4\tabcolsep) * \real{0.3333}}
  >{\RaggedRight\arraybackslash}p{(\linewidth - 4\tabcolsep) * \real{0.3333}}
  >{\RaggedRight\arraybackslash}p{(\linewidth - 4\tabcolsep) * \real{0.3333}}@{}}
\toprule\noalign{}
\begin{minipage}[b]{\linewidth}\RaggedRight
Neighbor count
\end{minipage} & \begin{minipage}[b]{\linewidth}\RaggedRight
Next state if dead
\end{minipage} & \begin{minipage}[b]{\linewidth}\RaggedRight
Next state if alive
\end{minipage} \\
\midrule\noalign{}
\endhead
\bottomrule\noalign{}
\endlastfoot
0 & 0 & 0 \\
1 & 0 & 0 \\
2 & 0 & 1 \\
3 & 1 & 1 \\
4 & 0 & 0 \\
5 & 0 & 0 \\
6 & 0 & 0 \\
7 & 0 & 0 \\
8 & 0 & 0 \\
\end{longtable}
}

The center is not included in the count. The expression is next = (n=3)
OR (current\_alive AND n=2). Current state and every neighbor belong to
the same old generation.
